\chapter{Arterial Blood Gas (ABG)}

\section*{What Is This Test?}
An arterial blood gas measures pH, partial pressure of oxygen, partial pressure of carbon dioxide, and bicarbonate-related acid-base information in arterial blood. It may also report oxygen saturation, base excess, lactate, and electrolytes. The test assesses ventilation, oxygenation, and acid-base balance.

\section*{Alternative Names}
ABG; Arterial Blood Gas Analysis; Blood Gas; Acid-Base Analysis.

\section*{Why This Test Is Ordered}
ABG testing is used for severe breathlessness, respiratory failure, shock, altered mental status, suspected metabolic crisis, sepsis, poisoning, and monitoring of ventilatory support. It is especially useful when a pulse-oximeter reading does not explain the patient's condition or when carbon dioxide retention is suspected.

\section*{How This Test Is Performed}
A small arterial blood sample is collected, commonly from the radial artery, into a heparinized blood-gas syringe and analyzed promptly. The laboratory instrument measures pH, oxygen, and carbon dioxide and calculates related values. Oxygen delivery, temperature, ventilation, and the patient's clinical state should be recorded.

\section*{Preparation, Risks, and Limitations}
No fasting is required. The sample must be free of air bubbles and analyzed without undue delay. Temporary pain, bruising, bleeding, or rarely arterial spasm may occur. Local anesthetic can reduce discomfort when appropriate. Supplemental oxygen and recent changes in ventilation substantially affect the result.

\section*{Understanding Results}
Low pH indicates acidemia and high pH alkalemia. Carbon dioxide primarily reflects respiratory influence, while bicarbonate and base excess reflect the metabolic component. Oxygen and oxygen saturation assess oxygenation but must be interpreted with inspired oxygen and the clinical context. A mixed disorder may be present even when pH is near normal.

\section*{References}
American Thoracic Society guidance on arterial blood gas interpretation and respiratory failure.

\clearpage
\section*{Understanding Results and Follow-up}
The laboratory report should identify the specimen, method, units, reference interval or decision limit, and whether the result is positive, negative, abnormal, or inconclusive. Reference intervals vary with age, sex, pregnancy, specimen type, assay, and laboratory; a result outside a range is not automatically a diagnosis, and a result within a range does not always exclude disease. Discuss unexpected results with the ordering clinician, who can decide whether repeat or confirmatory testing, treatment, or referral is appropriate. Do not change medicines or delay urgent care based on an isolated result.


\section*{Regulatory Status and Authorization Date}
Regulatory status applies to the specific assay, instrument, manufacturer, and intended use rather than to the test name alone. No single approval date can be assigned to this test category. For a commercial assay, verify the current product in the relevant regulator's database: the U.S. Food and Drug Administration (FDA) clearance, approval, or authorization; India's Central Drugs Standard Control Organisation (CDSCO) licence or permission; the European Union In Vitro Diagnostic Regulation (EU IVDR) CE certificate issued through the applicable conformity-assessment route; or the United Kingdom Medicines and Healthcare products Regulatory Agency (MHRA) registration and UKCA/CE status. Laboratory-developed tests may not have an individual FDA, CDSCO, EU IVDR, or MHRA approval date.
\clearpage
